\documentclass[runningheads]{llncs}

\usepackage[T1]{fontenc}
\usepackage{graphicx}
\usepackage{booktabs}
\usepackage{array}

\makeatletter
\renewenvironment{table}
               {\setlength\abovecaptionskip{10\p@}%
                \setlength\belowcaptionskip{0\p@}%
                \@float{table}}
               {\end@float}
\renewenvironment{table*}
               {\setlength\abovecaptionskip{10\p@}%
                \setlength\belowcaptionskip{0\p@}%
                \@dblfloat{table}}
               {\end@dblfloat}
\makeatother

\begin{document}

\title{Supplementary Material: RGB--Depth Generation for Cutaneous Neurofibroma}
\titlerunning{Supplementary Material}

\author{Anonymized Author(s)}
\authorrunning{Anonymized Author et al.}
\institute{Anonymized Affiliation(s)}

\maketitle

\appendix

\section{Model Settings}

All neural-network experiments used the training settings summarized below.

\begin{table}[!htbp]
\centering
\small
\renewcommand{\arraystretch}{1.08}
\begin{tabular}{>{\raggedright\arraybackslash}p{0.42\textwidth}>{\raggedright\arraybackslash}p{0.46\textwidth}}
\toprule
Setting & Value \\
\midrule
Input & RGB--depth crop \\
Split & 151/19/19 train/validation/test \\
Training resolution & $512 \times 512$ \\
Latent channels & 4 \\
Base channels & 32 \\
Parameters & 9,728,332 \\
Objective & RGB/depth reconstruction plus KL \\
Loss weights & RGB 1, depth 1, MSE 0.25, KL $10^{-4}$ \\
Optimizer & AdamW \\
Learning rate & $2\times10^{-4}$ \\
Weight decay & $10^{-4}$ \\
Batch size & 8 \\
Training schedule & 100 epochs \\
Gradient clipping & 1.0 \\
Precision & Mixed precision \\
Seed & 20260624 \\
Checkpoint selection & Validation loss \\
\bottomrule
\end{tabular}
\caption{Training settings for the four-channel VAE.}
\label{tab:vae-settings}
\end{table}

\begin{table}[!htbp]
\centering
\small
\renewcommand{\arraystretch}{1.08}
\begin{tabular}{>{\raggedright\arraybackslash}p{0.42\textwidth}>{\raggedright\arraybackslash}p{0.46\textwidth}}
\toprule
Setting & Value \\
\midrule
Input & RGB--depth crop plus guide mask \\
Split & Morphology-specific splits in main Table 2 \\
Training resolution & $256 \times 256$ \\
Latent channels & 4 \\
Mask channels & 4 \\
Base channels & 64 \\
Residual blocks per level & 2 \\
Time embedding & 256 dimensions \\
Dropout & 0.05 \\
Parameters & 12,653,316 \\
Objective & Masked source-to-target RGB/depth loss \\
Loss weights & Masked 1, unmasked 0, RGB 1, depth 1, latent 0 \\
Outside mask & Source preserved outside mask \\
Optimizer & AdamW \\
Learning rate & $10^{-4}$ \\
Weight decay & $10^{-4}$ \\
Batch size & 4 \\
Training schedule & 30 epochs \\
Gradient clipping & 1.0 \\
Precision & Mixed precision \\
Seed & 20260624 \\
Checkpoint selection & Validation loss \\
\bottomrule
\end{tabular}
\caption{Training settings for mask-guided inpainting.}
\label{tab:inpainting-settings}
\end{table}

\begin{table}[!htbp]
\centering
\small
\renewcommand{\arraystretch}{1.08}
\begin{tabular}{ll}
\toprule
Setting & Value \\
\midrule
Architecture & Three-level U-Net with GroupNorm and SiLU blocks \\
Input channels & 3 for RGB, 4 for RGB--D \\
Input resolution & $128 \times 128$ \\
Base channels & 24 \\
Loss & Weighted binary cross-entropy with logits plus Dice loss \\
Positive-weight cap & 50.0 \\
Optimizer & AdamW \\
Learning rate & $3\times10^{-4}$ \\
Weight decay & $10^{-4}$ \\
Batch size & 32 \\
Training schedule & 35 epochs \\
Decision threshold & 0.5 \\
Validation set & 53 held-out real crops \\
Synthetic augmentation & 512 crops, 128 per morphology \\
Seed & 20260628 \\
\bottomrule
\end{tabular}
\caption{Downstream detector settings used for the crop-level ablation.}
\label{tab:detectorsettings}
\end{table}

\clearpage

\section{Supplementary Figures}

\begin{figure}[!htbp]
\centering
\includegraphics[width=\textwidth]{figures/supp_model_workflows.png}
\caption{Model workflow schematics. The VAE reconstructs paired RGB and depth maps from a four-channel input, while the mask-guided inpainting runs use a source RGB--depth crop and guide mask to predict an RGB--depth lesion inside the requested region. Morphology is handled by separate class-specific runs.}
\label{fig:supp-model-workflows}
\end{figure}

\clearpage

\begin{figure}[!htbp]
\centering
\includegraphics[width=0.92\textwidth]{figures/supp_raycast_examples.png}
\caption{Lesion-centered examples from the geometry-controlled raycast branch. Rows show one rendered target for each morphology with its binary lesion mask and render-depth map. Depth is contrast-normalized within each crop for visualization, with light values displayed as closer.}
\label{fig:supp-raycast-examples}
\end{figure}

\clearpage

\begin{figure}[!htbp]
\centering
\includegraphics[width=0.92\textwidth]{figures/vae_test_reconstruction_progress_clean.png}
\caption{Held-out VAE reconstruction progression for one representative cNF crop. Columns show test RGB, reconstructed RGB, test pseudo-depth, and reconstructed pseudo-depth at selected checkpoints.}
\label{fig:supp-vae-progress}
\end{figure}

\clearpage

\begin{figure}[!htbp]
\centering
\includegraphics[width=\textwidth]{figures/downstream_ablation_training_curves.png}
\caption{Training loss and validation Dice curves for the downstream detector ablation. All four detector conditions used the same split, architecture family, threshold, and training schedule.}
\label{fig:supp-downstream-curves}
\end{figure}

\begin{table}[!htbp]
\centering
\small
\renewcommand{\arraystretch}{1.08}
\begin{tabular}{lrrrrrr}
\toprule
 & \multicolumn{3}{c}{RGB MAE} & \multicolumn{3}{c}{Depth MAE} \\
\cmidrule(lr){2-4}\cmidrule(lr){5-7}
Morphology & Source & Latent & Direct & Source & Latent & Direct \\
\midrule
Flat & 0.0509 & 0.1368 & 0.0511 & 0.0133 & 0.1214 & 0.0152 \\
Sessile & 0.0830 & 0.1388 & 0.0939 & 0.0905 & 0.1379 & 0.0895 \\
Globular & 0.0892 & 0.1241 & 0.1024 & 0.0735 & 0.2096 & 0.0735 \\
Pedunculated & 0.1341 & 0.1518 & 0.1238 & 0.1806 & 0.2544 & 0.1748 \\
\bottomrule
\end{tabular}
\caption{Per-morphology geometry-controlled validation on body-scan renderings. Each row uses 18 held-out examples; metrics are averaged inside the raycast lesion mask and lower values are better.}
\label{tab:supp-hsr-morphology}
\end{table}

\end{document}
